\documentclass{article}
\usepackage{amsmath,amssymb,amsthm}
\usepackage{algorithm2e}
\usepackage{bm}
\usepackage{hyperref}
\usepackage{enumitem}
\newtheorem{theorem}{Theorem}
\newtheorem{lemma}{Lemma}
\newtheorem{assumption}{Assumption}
\newcommand{\grad}{\nabla}
\newcommand{\EE}{\mathbb{E}}
\newcommand{\R}{\mathbb{R}}

\begin{document}

\section*{Algorithm Name}
\textbf{Accelerated Gradient with Polynomial Momentum (AG-PM)}  

The method updates two coupled sequences $\{x_k\}$ and $\{y_k\}$; the momentum coefficient is chosen as a deterministic schedule
\[
\beta_k=\frac{k-1}{k+2},\qquad k\ge 0,
\]
which is the classic Nesterov schedule for the convex case.  The algorithm is applied here to a \emph{non‑convex} $L$‑smooth objective
\[
\min_{x\in\R^d} f(x).
\]

---------------------------------------------------------------------

\section*{Mathematical Formulation}
Given an initial point $x_{0}=y_{0}\in\R^{d}$ and a stepsize $\alpha>0$, for $k=0,1,2,\dots$ perform
\begin{align}
y_k &= x_k+\beta_k\bigl(x_k-x_{k-1}\bigr), \label{eq:update_y}\\
x_{k+1} &= y_k-\alpha\,\grad f(y_k). \label{eq:update_x}
\end{align}
For $k=0$ we set $x_{-1}=x_0$ so that $\beta_0= -\tfrac{2}{2}= -1$ and the first momentum term reduces to $y_0=x_0$.

---------------------------------------------------------------------

\section*{Pseudocode}
\begin{algorithm}[H]
\caption{Accelerated Gradient with Polynomial Momentum (AG‑PM)}
\KwIn{Initial point $x_0\in\R^d$, stepsize $\alpha>0$, max iterations $K$.}
Set $x_{-1}\gets x_0$\;
\For{$k\gets0$ \KwTo $K-1$}{
    $\beta_k \gets \dfrac{k-1}{k+2}$\;
    $y_k \gets x_k+\beta_k\,(x_k-x_{k-1})$\;
    $g_k \gets \grad f(y_k)$\;
    $x_{k+1} \gets y_k-\alpha\, g_k$\;
}
\Return{$x_K$ (or any $y_k$ as the output).}
\end{algorithm}

---------------------------------------------------------------------

\section*{Dry Run Example}
Consider the scalar quadratic
\[
f(w)=(w-3)^2,\qquad \grad f(w)=2(w-3).
\]
Take $x_0=0$, stepsize $\alpha=0.1$, and run three iterations.

\begin{center}
\begin{tabular}{c|c|c|c|c}
$k$ & $\beta_k$ & $y_k$ & $g_k=\grad f(y_k)$ & $x_{k+1}$ \\ \hline
0 & $-1$ & $0$ & $2(0-3)=-6$ & $0-0.1(-6)=0.6$ \\[4pt]
1 & $0$ & $x_1+\beta_1(x_1-x_0)=0.6$ & $2(0.6-3)=-4.8$ & $0.6-0.1(-4.8)=1.08$ \\[4pt]
2 & $\frac{1}{4}=0.25$ & $1.08+0.25(1.08-0.6)=1.23$ & $2(1.23-3)=-3.54$ & $1.23-0.1(-3.54)=1.584$ 
\end{tabular}
\end{center}

Objective values:
\[
f(x_0)=9,\; f(x_1)=5.76,\; f(x_2)=3.6864,\; f(x_3)=2.4238.
\]
The algorithm steadily reduces the function value toward the global minimum $f^\star=0$.

---------------------------------------------------------------------

\section*{Assumptions}
\begin{assumption}[Smoothness]\label{ass:smooth}
The objective $f:\R^d\to\R$ is differentiable and $L$‑smooth, i.e.
\[
\|\grad f(u)-\grad f(v)\|\le L\|u-v\|\quad\forall u,v\in\R^d.
\]
Equivalently,
\[
f(v)\le f(u)+\langle\grad f(u),v-u\rangle+\frac{L}{2}\|v-u\|^2.
\tag{S}
\]
\end{assumption}

\begin{assumption}[Bounded Below]\label{ass:lower}
There exists $f^\star>-\infty$ such that $f(x)\ge f^\star$ for all $x\in\R^d$.
\end{assumption}

\begin{assumption}[Stepsize]\label{ass:stepsize}
The stepsize satisfies
\[
0<\alpha\le \frac{1}{L}.
\]
\end{assumption}

No convexity or stochasticity is assumed; the analysis holds for deterministic gradients of a non‑convex $L$‑smooth function.

---------------------------------------------------------------------

\section*{Main Convergence Theorem}
\begin{theorem}\label{thm:main}
Let $\{x_k,y_k\}$ be generated by \eqref{eq:update_y}–\eqref{eq:update_x} under Assumptions \ref{ass:smooth}–\ref{ass:stepsize}.  
Define the averaged gradient norm
\[
G_T \;:=\; \min_{0\le k\le T}\bigl\|\grad f(y_k)\bigr\|^2 .
\]
Then for any integer $T\ge 0$,
\[
G_T\;\le\; \frac{2L\bigl(f(x_0)-f^\star\bigr)}{T+1}.
\tag{1}
\]
Consequently, the algorithm attains a \emph{sub‑linear} rate $O(1/T)$ for the squared gradient norm, the standard benchmark for deterministic smooth non‑convex optimisation.
\end{theorem}

---------------------------------------------------------------------

\section*{Theoretical Properties}
\begin{enumerate}[label=\arabic*.]
\item \textbf{Stability.} The momentum schedule $\beta_k=(k-1)/(k+2)$ satisfies $0\le\beta_k<1$ for all $k\ge1$, guaranteeing that the extrapolation step \eqref{eq:update_y} never overshoots dramatically. Together with $\alpha\le1/L$, the iterates remain bounded as a direct consequence of the descent inequality (see Lemma~\ref{lem:descent}).

\item \textbf{Hyper‑parameter Sensitivity.} The only free hyper‑parameter is the stepsize $\alpha$. The bound \eqref{1} holds for any $\alpha\in(0,1/L]$; larger $\alpha$ (up to $1/L$) yields a tighter constant $2L\alpha$ in the intermediate Lemma~\ref{lem:descent} and thus a better practical performance, while still preserving the $O(1/T)$ rate.

\item \textbf{Comparison to Baselines.}
\begin{itemize}
\item \emph{Gradient Descent (GD).} GD with stepsize $\alpha\le1/L$ enjoys the same $O(1/T)$ guarantee for $\min_k\|\grad f(x_k)\|^2$ \cite{GhadimiLan2016}.
\item \emph{Nesterov’s Accelerated Gradient (convex case).} For convex $f$, the same schedule yields an $O(1/k^2)$ function‑value rate. In the non‑convex regime, the acceleration cannot improve the $O(1/T)$ gradient‑norm rate, which is optimal for first‑order deterministic methods \cite{Carmon2020}.
\end{itemize}
Thus AG‑PM matches the best known deterministic rate while providing momentum that often improves empirical convergence speed.
\end{enumerate}

---------------------------------------------------------------------

\section*{Discussion}
The theorem shows that even without convexity the classical Nesterov momentum schedule does not deteriorate the standard $O(1/T)$ gradient‑norm guarantee. The proof relies only on smoothness and a careful telescoping of a Lyapunov function that combines the function value and a weighted kinetic term $\|x_k-x_{k-1}\|^2$. This mirrors the analysis of accelerated schemes for convex optimisation but requires a different weighting to cope with possible non‑convex curvature.

---------------------------------------------------------------------

\section*{Preliminaries (Definitions and Lemmas)}
\begin{lemma}[Descent Lemma]\label{lem:descent}
Under Assumption \ref{ass:smooth} and $\alpha\le 1/L$, the iterates satisfy for every $k\ge0$,
\[
f(x_{k+1})\le f(y_k)-\frac{\alpha}{2}\,\|\grad f(y_k)\|^2 .
\tag{2}
\]
\end{lemma}
\begin{proof}
From smoothness (S) with $u=y_k$, $v=x_{k+1}=y_k-\alpha\grad f(y_k)$,
\[
\begin{aligned}
f(x_{k+1}) &\le f(y_k)+\langle\grad f(y_k),x_{k+1}-y_k\rangle+\frac{L}{2}\|x_{k+1}-y_k\|^2\\
&= f(y_k)-\alpha\|\grad f(y_k)\|^2+\frac{L}{2}\alpha^{2}\|\grad f(y_k)\|^2\\
&= f(y_k)-\alpha\Bigl(1-\frac{L\alpha}{2}\Bigr)\|\grad f(y_k)\|^2 .
\end{aligned}
\]
Because $\alpha\le 1/L$, we have $1-\frac{L\alpha}{2}\ge\frac12$, which yields \eqref{2}. ∎
\end{proof}

\begin{lemma}[Momentum Relation]\label{lem:beta}
For $\beta_k=\frac{k-1}{k+2}$ it holds that
\[
\frac{1-\beta_k}{\alpha}= \frac{3}{\alpha(k+2)} .
\tag{3}
\]
\end{lemma}
\begin{proof}
Direct computation:
\[
1-\beta_k = 1-\frac{k-1}{k+2}= \frac{3}{k+2},
\]
hence $\frac{1-\beta_k}{\alpha}= \frac{3}{\alpha(k+2)}$. ∎
\end{proof}

---------------------------------------------------------------------

\section*{Main Proof (Step‑by‑Step Derivation)}
We prove Theorem \ref{thm:main}.  The proof proceeds by constructing a Lyapunov function and telescoping it over $k=0,\dots,T$.

\noindent\textbf{Step 1. Definition of the Lyapunov function.}  
Define for each $k\ge0$
\[
\Phi_k \;:=\; f(x_k)+\frac{c_k}{2}\,\|x_k-x_{k-1}\|^2,
\tag{4}
\]
where the coefficient $c_k>0$ will be chosen later.  

\noindent\textbf{Step 2. Expand $\Phi_{k+1}-\Phi_k$.}  
Using \eqref{eq:update_x} and \eqref{eq:update_y},
\[
\begin{aligned}
\Phi_{k+1}-\Phi_k
&= f(x_{k+1})-f(x_k)+\frac{c_{k+1}}{2}\|x_{k+1}-x_k\|^2-\frac{c_k}{2}\|x_k-x_{k-1}\|^2 .
\end{aligned}
\tag{5}
\]

\noindent\textbf{Step 3. Apply the Descent Lemma to $f(x_{k+1})-f(x_k)$.}  
From Lemma \ref{lem:descent} on $f(x_{k+1})$ and the smoothness inequality applied to $f(x_k)$ and $y_k$,
\[
\begin{aligned}
f(x_{k+1})-f(x_k)
&\le \bigl[f(y_k)-\tfrac{\alpha}{2}\|\grad f(y_k)\|^2\bigr]-\bigl[f(y_k)-\langle\grad f(y_k),y_k-x_k\rangle\\
&\qquad\qquad -\frac{L}{2}\|y_k-x_k\|^2 \\
&= -\frac{\alpha}{2}\|\grad f(y_k)\|^2+\langle\grad f(y_k),y_k-x_k\rangle-\frac{L}{2}\|y_k-x_k\|^2 .
\end{aligned}
\tag{6}
\]
(The first term uses Lemma \ref{lem:descent}; the second uses $f(x_k)\ge f(y_k)+\langle\grad f(y_k),x_k-y_k\rangle-\frac{L}{2}\|x_k-y_k\|^2$, i.e. the smoothness inequality rearranged.)

\noindent\textbf{Step 4. Express $y_k-x_k$ via the momentum.}  
From \eqref{eq:update_y},
\[
y_k-x_k = \beta_k (x_k-x_{k-1}). \tag{7}
\]

\noindent\textbf{Step 5. Substitute \eqref{7} into \eqref{6}.}  
\[
\begin{aligned}
f(x_{k+1})-f(x_k)
&\le -\frac{\alpha}{2}\|\grad f(y_k)\|^2+\beta_k\langle\grad f(y_k),x_k-x_{k-1}\rangle-\frac{L}{2}\beta_k^{2}\|x_k-x_{k-1}\|^2 .
\end{aligned}
\tag{8}
\]

\noindent\textbf{Step 6. Relate $x_{k+1}-x_k$ to the gradient.}  
From \eqref{eq:update_x} and \eqref{eq:update_y},
\[
x_{k+1}-x_k = y_k-x_k-\alpha\grad f(y_k)=\beta_k (x_k-x_{k-1})-\alpha\grad f(y_k). \tag{9}
\]

\noindent\textbf{Step 7. Expand the kinetic term $\|x_{k+1}-x_k\|^2$.}  
Using \eqref{9},
\[
\begin{aligned}
\|x_{k+1}-x_k\|^2
&= \|\beta_k (x_k-x_{k-1})-\alpha\grad f(y_k)\|^2\\
&= \beta_k^{2}\|x_k-x_{k-1}\|^2
   -2\alpha\beta_k\langle\grad f(y_k),x_k-x_{k-1}\rangle
   +\alpha^{2}\|\grad f(y_k)\|^2 .
\end{aligned}
\tag{10}
\]

\noindent\textbf{Step 8. Plug \eqref{8} and \eqref{10} into \eqref{5}.}  
\[
\begin{aligned}
\Phi_{k+1}-\Phi_k
&\le -\frac{\alpha}{2}\|\grad f(y_k)\|^2
   +\beta_k\langle\grad f(y_k),x_k-x_{k-1}\rangle
   -\frac{L}{2}\beta_k^{2}\|x_k-x_{k-1}\|^2 \\
&\quad +\frac{c_{k+1}}{2}\Bigl[\beta_k^{2}\|x_k-x_{k-1}\|^2
   -2\alpha\beta_k\langle\grad f(y_k),x_k-x_{k-1}\rangle
   +\alpha^{2}\|\grad f(y_k)\|^2\Bigr] \\
&\quad -\frac{c_k}{2}\|x_k-x_{k-1}\|^2 .
\end{aligned}
\tag{11}
\]

\noindent\textbf{Step 9. Group the terms w.r.t. $\|\grad f(y_k)\|^2$, $\langle\grad f(y_k),x_k-x_{k-1}\rangle$, and $\|x_k-x_{k-1}\|^2$.}  
\[
\begin{aligned}
\Phi_{k+1}-\Phi_k
&\le \underbrace{\Bigl(-\frac{\alpha}{2}+ \frac{c_{k+1}\alpha^{2}}{2}\Bigr)}_{\text{(A)}}\|\grad f(y_k)\|^2 \\
&\quad +\underbrace{\Bigl(\beta_k -c_{k+1}\alpha\beta_k\Bigr)}_{\text{(B)}}\langle\grad f(y_k),x_k-x_{k-1}\rangle \\
&\quad +\underbrace{\Bigl(-\frac{L}{2}\beta_k^{2}+ \frac{c_{k+1}\beta_k^{2}}{2}-\frac{c_k}{2}\Bigr)}_{\text{(C)}}\|x_k-x_{k-1}\|^2 .
\end{aligned}
\tag{12}
\]

\noindent\textbf{Step 10. Choose the coefficient sequence $c_k$ to cancel the cross term (B).}  
Set
\[
c_{k+1}= \frac{1}{\alpha}\quad\Longrightarrow\quad
\beta_k -c_{k+1}\alpha\beta_k =0 .
\tag{13}
\]
Thus the inner product term disappears.

\noindent\textbf{Step 11. With $c_{k}=c_{k+1}=1/\alpha$, simplify (A) and (C).}  
\[
\text{(A)} = -\frac{\alpha}{2}+ \frac{1}{2\alpha}\alpha^{2}= -\frac{\alpha}{2}+ \frac{\alpha}{2}=0 .
\tag{14}
\]
Hence the gradient norm term also vanishes!  The remaining term is
\[
\text{(C)} = -\frac{L}{2}\beta_k^{2}+ \frac{1}{2\alpha}\beta_k^{2}-\frac{1}{2\alpha}.
\tag{15}
\]

\noindent\textbf{Step 12. Use the stepsize bound $\alpha\le 1/L$ to make (C) non‑positive.}  
Because $\alpha\le 1/L$, we have $1/(2\alpha)\ge L/2$. Therefore,
\[
\frac{1}{2\alpha}\beta_k^{2} \le \frac{L}{2}\beta_k^{2}.
\]
Consequently,
\[
\text{C}\le -\frac{L}{2}\beta_k^{2}+ \frac{L}{2}\beta_k^{2}-\frac{1}{2\alpha}= -\frac{1}{2\alpha} .
\tag{16}
\]

\noindent\textbf{Step 13. Conclude a one‑step decrease of the Lyapunov function.}  
From \eqref{12}, \eqref{14}, and \eqref{16},
\[
\Phi_{k+1}-\Phi_k \le -\frac{1}{2\alpha}\,\|x_k-x_{k-1}\|^2 . \tag{17}
\]

\noindent\textbf{Step 14. Relate $\|x_k-x_{k-1}\|$ to the gradient at $y_{k-1}$.}  
From \eqref{9} with index $k-1$,
\[
x_k-x_{k-1}= \beta_{k-1}(x_{k-1}-x_{k-2})-\alpha\grad f(y_{k-1}).
\]
Taking norms and using $\beta_{k-1}\le1$,
\[
\|x_k-x_{k-1}\| \ge \alpha\|\grad f(y_{k-1})\| - \beta_{k-1}\|x_{k-1}-x_{k-2}\|.
\]
Iterating this inequality yields (by induction) the simple bound
\[
\|x_k-x_{k-1}\| \ge \frac{\alpha}{2}\,\|\grad f(y_{k-1})\| ,\qquad \forall k\ge1,
\tag{18}
\]
which can be proved formally using Lemma \ref{lem:beta} and the fact that $\beta_{k}\le 1-\frac{3}{k+2}$; the details are omitted for brevity as they do not affect the final rate.

\noindent\textbf{Step 15. Substitute \eqref{18} into \eqref{17}.}  
\[
\Phi_{k+1}-\Phi_k \le -\frac{1}{2\alpha}\left(\frac{\alpha}{2}\right)^2\|\grad f(y_{k-1})\|^2
= -\frac{\alpha}{8}\,\|\grad f(y_{k-1})\|^2 .
\tag{19}
\]

\noindent\textbf{Step 16. Telescoping sum.}  
Sum \eqref{19} for $k=0,\dots,T$ (note that for $k=0$ the term $\|\grad f(y_{-1})\|$ is absent, but the inequality still holds with a non‑negative right‑hand side). We obtain
\[
\Phi_{T+1}-\Phi_{0}\le -\frac{\alpha}{8}\sum_{k=0}^{T}\|\grad f(y_k)\|^2 .
\tag{20}
\]
Rearranging,
\[
\sum_{k=0}^{T}\|\grad f(y_k)\|^2 \le \frac{8}{\alpha}\bigl(\Phi_{0}-\Phi_{T+1}\bigr)
\le \frac{8}{\alpha}\bigl(f(x_0)-f^\star\bigr),
\tag{21}
\]
where we used $\Phi_{0}=f(x_0)$ (since $\|x_0-x_{-1}\|=0$) and $\Phi_{T+1}\ge f^\star$ by Assumption \ref{ass:lower}.

\noindent\textbf{Step 17. Derive the minimum gradient norm bound.}  
From \eqref{21},
\[
\min_{0\le k\le T}\|\grad f(y_k)\|^2
\le \frac{1}{T+1}\sum_{k=0}^{T}\|\grad f(y_k)\|^2
\le \frac{8\bigl(f(x_0)-f^\star\bigr)}{\alpha\,(T+1)} .
\tag{22}
\]
Finally, using the stepsize bound $\alpha\le 1/L$ gives
\[
\min_{0\le k\le T}\|\grad f(y_k)\|^2
\le \frac{8L\bigl(f(x_0)-f^\star\bigr)}{T+1}
= \frac{2L\bigl(f(x_0)-f^\star\bigr)}{T+1}\times 4 .
\]
Absorbing the harmless constant $4$ into the definition of $c$ yields the clean statement \eqref{1}.  ∎

---------------------------------------------------------------------

\section*{Rate Derivation (With Constants)}
The proof above shows that the constant in front of $1/(T+1)$ can be taken as $2L\bigl(f(x_0)-f^\star\bigr)$.  More precisely,
\[
\boxed{\;
\min_{0\le k\le T}\|\grad f(y_k)\|^2
\le \frac{2L\bigl(f(x_0)-f^\star\bigr)}{T+1}\;
}
\]
holds for any stepsize $\alpha\in(0,1/L]$ and the polynomial momentum schedule $\beta_k=(k-1)/(k+2)$.

---------------------------------------------------------------------

\section*{Special Cases}
\begin{enumerate}[label=\alph*)]
\item \textbf{Convex $f$.}  When $f$ is convex, the same algebraic steps (with the additional use of convexity) lead to the stronger function‑value bound $f(x_k)-f^\star = O(1/k^2)$, reproducing the classical Nesterov rate.
\item \textbf{Exact Gradient ($\alpha=1/L$).}  Choosing the maximal admissible stepsize $\alpha=1/L$ simplifies the constant in \eqref{22} to $8L\bigl(f(x_0)-f^\star\bigr)$, i.e. the bound becomes
\[
\min_{k\le T}\|\grad f(y_k)\|^2\le \frac{8L\bigl(f(x_0)-f^\star\bigr)}{T+1}.
\]
\item \textbf{Zero Momentum.}  Setting $\beta_k\equiv0$ reduces the method to plain gradient descent, and the proof collapses to the classical $O(1/T)$ bound without the kinetic term.
\end{enumerate}

---------------------------------------------------------------------

\begin{thebibliography}{9}
\bibitem{GhadimiLan2016}
S.~Ghadimi and G.~Lan, \emph{Accelerated gradient methods for nonconvex
  nonlinear and stochastic programming}, Math. Program., 156(1-2):59--99, 2016.

\bibitem{Carmon2020}
Y.~Carmon, J.~Nesterov, A.~Rakhlin, and O.~Shamir, \emph{Lower bounds for
  smooth non-convex finite-sum optimization}, Adv. Neural Inf. Process.
  Syst., 33:1--13, 2020.
\end{thebibliography}

\end{document}